\begin{frame}{자주 묻는 질문}
  \begin{itemize}
    \item \textbf{Q. 두 모델 성능이 거의 같으면?} 오히려 좋은 소재
      --- Week 6.3: 차이가 검증셋 표본 크기에 비해 작다면
      ``우연''일 수도 있다는 것 자체가 \textbf{통계적으로 중요한
      관찰}. 예측이 어떤 포인트에서 갈리는지, 어느 쪽이 해석하기
      쉬운지(로지스틱의 계수 vs GBDT의 SHAP)로 채운다.
    \item \textbf{Q. 결측치·이상치가 많은 데이터는 감점?} 아니다
      --- 오히려 전처리 결정(어떻게 채웠는지, 이상치를 제거·유지
      했는지)과 \textbf{그 근거}가 좋은 평가 포인트. ``지저분한
      데이터를 다룬 판단 과정''이 실전에 훨씬 가깝다.
    \item \textbf{Q. 역할 분담은?} 데이터·코드·보고서·발표를
      \textbf{팀 전체가 함께 결정}하되 작성은 나눠도 된다.
      ``코드는 A, 보고서는 B''로 완전히 분리하면 B가 A의 코드를
      모르는 채 정당화를 쓰게 됨(=``나쁜 설명'') --- M2
      중간코드 리뷰가 이를 강제.
    \item \textbf{Q. 4주 안에 안 끝날 것 같은데?} 성공 기준은
      ``높은 성능''이 아니라 \textbf{``정직한 절차''} --- 모델을
      2개로 줄이고(최소선은 2개) ``시간 부족으로 3번째 모델을
      못 돌렸다''를 정직하게 쓰는 것이 부풀리기보다 낫다.
  \end{itemize}
\end{frame}
